\documentclass[11pt]{article}
\usepackage[margin=0.7in]{geometry}
\usepackage{titlesec}
\usepackage{enumitem}
\usepackage{parskip}

\titleformat{\section}{\Large\bfseries}{\thesection}{1em}{}
\titleformat{\subsection}{\large\bfseries}{\thesubsection}{1em}{}

\setlength{\parskip}{0.3em}

\title{\textbf{Week 6: Modal Jazz (Late 1950s-1960s)} \\ \large Quick Reference}
\author{}
\date{}

\begin{document}

\maketitle

\section*{Period Overview}

\textbf{What is Modal Jazz?} Freedom through simplicity. Uses modes (scales) instead of rapid chord changes. Static harmony allows melodic exploration - musicians no longer "chasing changes." Spacious, floating, contemplative feel. Spiritual seeking, interest in Eastern philosophy. Key albums: \emph{Kind of Blue} (Miles, 1959), \emph{My Favorite Things} (Coltrane, 1960).

\textbf{Key Sound:} Spacious, floating, meditative. Minimal chord changes, one mode for extended periods, open atmosphere.

\textbf{Relation to Bebop:} SIMPLIFICATION - replaces bebop's complex rapid chord changes with static modes. Freed improvisers from "chasing changes" to develop melodies. Less dense, more spacious than bebop. Contemplative vs bebop's intensity.

\section*{Required Listening (11 Pieces)}

\subsection*{Miles Davis - "So What" from \emph{Kind of Blue} (1959)}
This modal jazz masterpiece represents a revolutionary departure from bebop through its use of static harmony with only two modes across the entire 32-bar AABA form (16-8-8 bars). The instrumentation features Miles' distinctive Harmon mute, Bill Evans' impressionistic piano voicings influenced by Ravel and Debussy, and a two-note bass ostinato that establishes the D Dorian mode. The relaxed spacious pulse and cool contemplative atmosphere create a floating meditative quality characteristic of modal jazz. Rather than "chasing changes" through rapid bebop chord progressions, improvisers develop horizontal melodic lines over static harmony, freed from bebop's vertical chord outlining. This piece from the 1959 "multiple revolutions" year demonstrates how modal simplification allowed contemplative spiritual exploration rather than bebop's frenetic intellectual intensity.

\subsection*{Miles Davis - "Riot" from \emph{Nefertiti} (1967)}
This modal jazz piece demonstrates the continued evolution of Miles' modal approach into the late 1960s with a more aggressive edge. The instrumentation features the second great quintet with Wayne Shorter's tenor saxophone, Herbie Hancock's piano, Ron Carter's bass, and Tony Williams' explosive drumming. The modal structure maintains static harmony over extended periods, but the rhythmic interplay is far more intense and fragmented than earlier modal works. This piece shows modal jazz's development from the spacious contemplative feel of "Kind of Blue" toward greater rhythmic complexity and tension while still avoiding bebop's rapid chord changes. The title "Riot" reflects both the musical intensity and the mid-1960s Civil Rights era context, where jazz's evolution paralleled social liberation movements.

\subsection*{John Coltrane - "My Favorite Things" (1960)}
This modal jazz piece features Coltrane on soprano saxophone playing in 3/4 waltz time over an E minor/major modal vamp, creating a hypnotic repetitive quality. McCoy Tyner's piano uses quartal harmony built on intervals of fourths rather than bebop's tertian chord structures, contributing to the ethereal mystical atmosphere. The extended solos lasting over 10 minutes demonstrate how modal jazz allowed spiritual exploration over static harmony rather than bebop's head-solo-head economy. Coltrane abandoned the extreme harmonic complexity of "Giant Steps" for this modal simplicity, showing his journey from bebop's ultimate chord-based improvisation toward spiritual seeking. This Rodgers and Hammerstein melody developed endlessly over a vamp represents modal jazz's interest in Eastern philosophy and meditation through repetition.

\subsection*{Bill Evans - "Nardis" from \emph{Explorations} (1961)}
This modal jazz piano trio piece exemplifies conversational interaction between equal voices rather than bebop's soloist-accompanist hierarchy. Evans' impressionistic delicate piano touch shows influence from classical composers Ravel and Debussy rather than bebop's blues roots. The instrumentation features melodic bass playing rather than strict walking bass patterns and subtle brushwork on drums, creating a sparse intimate chamber-music texture. The modal composition eliminates bebop's rapid chord changes in favor of static harmony that allows spacious melodic exploration. This piece demonstrates how modal jazz replaced bebop's dense rapid interplay with conversational intimacy and impressionistic color exploration.

\subsection*{Dave Brubeck - "Take Five" from \emph{Time Out} (1959)}
This modal jazz piece in 5/4 time signature demonstrates cool jazz's West Coast approach to modal concepts with accessible odd meter experimentation. Paul Desmond's light smooth alto saxophone sound and the strong bass ostinato create a medium tempo vamp structure that's accessible despite the unusual meter. The largely modal harmonic approach with static harmony replaces bebop's complex rapid chord changes characteristic of East Coast bebop. This piece shows the cool jazz aesthetic of pleasant laid-back feeling versus bebop's intensity, making experimental concepts commercially successful. The 1959 recording represents one of the "multiple revolutions" in jazz that year, bringing modal concepts and odd meters to mainstream audiences.

\subsection*{Herbie Hancock - "Maiden Voyage" (1965)}
This modal jazz composition features suspended chords and quartal harmony built on fourths, creating a floating spacious feel with harmonic ambiguity. The ocean imagery and wave-like rhythm at medium tempo demonstrate modal jazz's exploration of sound color and atmosphere over bebop's virtuosic line construction. The instrumentation creates an open airy ensemble texture that replaces bebop's functional chord progressions and dense harmonic movement. The static ambiguous harmony freed improvisers from bebop's clear chord changes, allowing melodic and textural exploration. This piece exemplifies mid-1960s modal jazz's maturity in using suspended and quartal voicings to create contemplative spacious atmospheres.

\subsection*{Stan Getz - "The Girl from Ipanema" (1964)}
This bossa nova piece fuses Brazilian music with cool jazz modal aesthetics, creating the opposite of bebop's intensity through gentle accessible grooves. Getz's breathy relaxed tenor saxophone timbre shows Lester Young's influence with nonchalant phrasing contrasting bebop's crisp articulation. The instrumentation features guitar-based textures and Astrud Gilberto's simple light vocal over a soft swaying groove with Antonio Carlos Jobim's simple melody. The harmonic approach uses modal-ish simple chord changes rather than bebop's complex progressions, emphasizing feel and space over technical virtuosity. This commercially successful fusion brought modal jazz concepts to mainstream audiences through Brazilian rhythms and accessible melodies.

\subsection*{Andrew Hill - "Refuge" from \emph{Point of Departure} (1964)}
This transitional piece pushes modal jazz toward free jazz through dense angular dissonant piano clusters and Eric Dolphy's presence. While still modal in structure, the piece is much harsher with rhythmic complexity that alternates freer sections with structured motifs. The dissonance exceeds bebop's complexity through impressionist and classical influences rather than bebop's blues roots. Hill's compositional approach uses collective improvisation elements that extend bebop's interactive interplay toward near-chaos while maintaining some modal structure. This 1964 recording demonstrates the boundary between modal jazz's static harmony and free jazz's complete abandonment of predetermined structures.

\subsection*{Joe Henderson - "Serenity" from \emph{In 'N Out} (1964)}
This modal jazz composition demonstrates the post-bop movement's balance between modal simplicity and harmonic sophistication. Henderson's tenor saxophone explores melodic development over static modal harmony rather than navigating bebop's rapid chord changes. The piece maintains modal jazz's spacious contemplative quality while showing more compositional complexity than simple vamp structures. The instrumentation allows extended improvisation over modal foundations, freeing soloists from "chasing changes" to develop melodic ideas. This 1964 recording shows modal jazz's maturity as a distinct approach separate from both bebop's harmonic complexity and free jazz's complete structural abandonment.

\subsection*{Wayne Shorter - "Speak No Evil" (1964)}
This modal jazz composition features mysterious film noir qualities and enigmatic dramatic motivic development rather than bebop's straightforward head-solos format. Shorter's tenor saxophone has similar timbre to Coltrane but with more singable rhythmic phrasing, accompanied by Elvin Jones' driving drums. The modal simplification of harmony replaces bebop's rapid chord changes with static structures that allow atmospheric and compositional exploration. Recorded approximately two weeks after Coltrane's "A Love Supreme," this piece shares the modal era aesthetic but demonstrates a different compositional approach. This shows how modal jazz allowed diverse artistic voices beyond bebop's more standardized formal structures.

\subsection*{John Coltrane - "Acknowledgment" from \emph{A Love Supreme} (1964)}
This modal jazz spiritual suite represents Coltrane's journey from bebop's "Giant Steps" to devotional modal exploration through a simple four-note modal motif. Coltrane's tenor saxophone features a thicker harder tone than his earlier bebop period with "upward reaching" phrasing into screaming high register. The rhythm section maintains a Latin/hard-bop groove tinge beneath the modal structure, while the famous double-tracked vocal chant "A Love Supreme" emphasizes spiritual devotion. The piece abandons bebop's complex chord-based improvisation for extended modal development of a simple motif, allowing contemplative spiritual seeking. This represents Coltrane's last major modal work before moving to free jazz, demonstrating modal jazz's role in his spiritual and musical evolution.

\section*{Key Terms}

\textbf{Mode} - Scale type used as harmonic foundation (Dorian most common) \\
\textbf{Dorian Mode} - Minor scale with raised 6th \\
\textbf{Vamp} - Repeated chord/rhythm pattern \\
\textbf{Ostinato} - Repeated musical phrase (often bass) \\
\textbf{Pedal Point} - Sustained bass note, harmony changes above \\
\textbf{Quartal Harmony} - Chords built on 4ths (McCoy Tyner) \\
\textbf{Suspended Chords} - Creates floating quality (Hancock) \\
\textbf{Bossa Nova} - Brazilian + jazz fusion, softer than samba

\section*{Readings}

\textbf{Szwed, \emph{Jazz 101} Ch. 22-23} - Post-bebop improvisation, 1959 revolutions \\
\textbf{Berliner, \emph{Thinking in Jazz} Ch. 3} - "A Very Structured Thing"

\section*{How Modal Jazz Emerged from Hard Bop}

Modal simplified Hard Bop's complex chord progressions. Musicians (Miles, Coltrane) sought melodic freedom and space rather than "chasing changes." Allowed contemplative spacious improvisation and spiritual exploration. Static harmony replaced rapid changes.

\end{document}
